\section*{Abstract}\label{abstract}

Die fortschreitende Verbreitung von Internet of Things (IoT)-Geräten erweitert die Angriffsfläche moderner Netzwerke erheblich \cite{hamedani2025iotsurvey}. Klassische Intrusion Detection Systeme (IDS) auf Basis von Signaturvergleich erreichen in diesem Kontext methodische Grenzen, weshalb Machine Learning (ML) als datengetriebene Alternative untersucht wird \cite{kikissagbe2024review}. Auf dem CICIoT2023-Datensatz \cite{neto2023ciciot2023} zeigen Raturi et al. \cite{raturi2026exploratory}, dass baumbasierte Modelle und Gradient Boosting im Speziellen die höchste Balanced Accuracy in der Mehrklassenklassifikation erzielen. Die Studie identifiziert jedoch drei methodische Lücken, darunter die fehlende Erklärbarkeit der Modellentscheidungen \cite{raturi2026exploratory}.

Die vorliegende Arbeit setzt an dieser Lücke an. Sie integriert SHapley Additive Explanations (SHAP) als post-hoc-Erklärungskomponente in einen Gradient-Boosting-Klassifikator auf dem CICIoT2023-Datensatz und bewertet die Modellleistung mit der Balanced Accuracy als primärer Metrik. Die methodische Vorlage für die SHAP-Integration stützt sich auf Mohale und Obagbuwa \cite{mohale2025xai}, die quantitative Bewertung der Erklärungsqualität auf Hermosilla, Berríos und Allende-Cid \cite{hermosilla2025xai_forensic}. Die Verbreitung und wissenschaftliche Begründung von SHAP in der IDS-Forschung wird durch den systematischen Review von Ogunseyi et al. \cite{ogunseyi2026xai_review} gestützt. Eine Ablation Study untersucht zusätzlich den Einfluss der Principal Component Analysis (PCA) auf Modellleistung und Erklärungen, da Raturi et al. \cite{raturi2026exploratory} PCA verwenden, ohne deren Wirkung systematisch zu analysieren. Vergleichswerte für Gradient Boosting und PCA auf dem CICIoT2023-Datensatz werden aus Alharby \cite{alharby2025ciciot} gewonnen. Der Umgang mit Klassenungleichgewicht wird im Diskussionsteil im Vergleich zu Almahaqeri et al. \cite{almahaqeri2026gradient} und unter Bezug auf Hosseini et al. \cite{hosseini2025imbalance} reflektiert.

\textbf{Schlüsselwörter:} Internet of Things, Machine Learning, Gradient Boosting, CICIoT2023, Balanced Accuracy, Intrusion Detection, Explainable Artificial Intelligence, SHAP, Principal Component Analysis, Cyberangriffserkennung
